\documentclass[main.tex]{subfiles}


\begin{document}

% \AddToShipoutPictureFG{%
% \begin{tikzpicture}[overlay,remember picture]
% \path (current page.south west) -- (current page.north east)
% node[midway,scale=1,color=gray,sloped,opacity=0.4] {\textnormal{Кравченко Игорь Игоревич\hspace{1cm} +79010144910\hspace{1cm} super.antig@yandex.ru}};
% \end{tikzpicture}
% }


% %from pagenumber to up
% \phantomsection 
% \hypertarget{toc}{}

 

 

 
   

\section{Силы тяготения и реакции}


\begin{law}
\textbf{Закон всемирного тяготения.} Два тела массами~$m_1$ и~$m_2$, расположенные на расстоянии~$r$, притягиваются
друг к другу с \textbf{силой тяготения}
\begin{equation}
    \label{26grav_e1}
    F_\text{т}=G\dfrac{m_1m_2}{r^2},
\end{equation}
где~$G$~--- \emph{гравитационная постоянная} (см.~справочные таблицы). 
\end{law}

На рис.\,\ref{fig:p35} условно изображены планета и космический корабль.

    \begin{figure}[H]
    \centering

    \begin{tikzpicture}[font=\small,            line cap=round,                             >={Stealth[inset=0pt,round,length=3mm, angle'=20]},vec/.style={>={Stealth[inset=0pt,round,length=3.5mm, angle'=20]}}]


\shade[ball color=Green] (0,0) circle (.6);
\node [left,black] at (-.6,0) {П};

\path (0,0) ++ (20:3cm) node (s) {} ++ (.15,0) node[black,right] {К};
\shade[ball color=lightgray] (s) circle (.15);

%r
\draw[densely dashed] (0,0) --++ (110:1.5cm);
\draw[densely dashed] (s.center) --++ (110:1.5cm);

\draw[<->] (0,0) ++ (110:1.3cm) --++ (20:3cm) node[midway,above] {$r$};


%vec
\node at (0,0) [circle,fill=red,inner sep=0pt,minimum size=1mm,label=below:] {};
\draw[->,red,thick,vec] (0,0) --++ (20:1cm) node[black,above] {$F_\text{т}$};
\node at (s) [circle,fill=red,inner sep=0pt,minimum size=1mm,label=below:] {};
\draw[->,red,thick,vec] (s.center) --++ (200:1cm) node[black,below] {$F_\text{т}$};

%m
\node [below,black] at (0,-.6) {$m_\text{п}$};
\path (0,0) ++ (20:3cm) ++ (0,-.15) node[black,below] {$m_\text{к}$};



    \end{tikzpicture}
\caption{Притяжение планеты и космического корабля}
\label{fig:p35}
\end{figure}

Планета~П массой~$m_\text{п}$ и корабль~К массой~$m_\text{к}$ расположены на некотором расстоянии~$r$ друг от друга\footnote{Для однородных шарообразных тел расстояние~$r$ есть расстояние между их центрами.}. Как и любые два~тела, обладающих массой, планета и корабль \emph{взаимодействуют} друг с другом так, что планета притягивает корабль с силой~$F_\text{т}$, а корабль~--- планету с такой же силой~$F_\text{т}$. 

При свободном падении у поверхности планеты все тела движутся с одинаковым ускорением~$g$. Тогда по второму закону~Ньютона
сила, действующая на тело массы~$m$  со стороны планеты и называемая \textbf{силой тяжести}, равна:
\begin{equation}
    \label{26grav_e2}
    F_\text{т}=mg.
\end{equation}

Вообще, сила тяжести и сила тяготения~--- это взаимозаменяемые термины, обозначающие одну и ту же силу \emph{гравитационного взаимодействия}. Так, силу тяжести, действующую на тело, можно называть и силой тяготения: численное значение этой силы гравитационного притяжения от этого не поменяется\footnote{Приравняв правые части формул~\eqref{26grav_e1} и~\eqref{26grav_e2}, можно получить формулу для ускорения свободного падения: $g=G\dfrac{m_\text{планеты}\strut}{r^2}$.}.

\sbox{0}{%
    \begin{tikzpicture}[font=\small,            line cap=round,                             >={Stealth[inset=0pt,round,length=3mm, angle'=20]},vec/.style={>={Stealth[inset=0pt,round,length=3.5mm, angle'=20]}}]


\filldraw[lightgray,semitransparent] (5,0) rectangle (8,-0.3);
\draw (5,0) -- (8,0);

\shade[ball color=lightgray,yshift=0] (6.5,.3) circle (.3);

%vec
\node at (6.5,-.2) [circle,fill=NavyBlue,inner sep=0pt,minimum size=1mm,label=below:] {};
\draw[->,NavyBlue,thick,vec] (6.5,-.2) --++ (0,-1) node[black,right] {$\vec P$};

\node at (6.5,.2) [circle,fill=red,inner sep=0pt,minimum size=1mm,label=below:] {};
\draw[->,red,thick,vec] (6.5,.2) --++ (0,1) node[black,right] {$\vec N$};


    \end{tikzpicture}%
}

{%
\setlength\intextsep{0pt}
\begin{wrapfigure}[]{r}{\wd0}
    \centering
    \usebox{0}%
    \caption{Шар\\ на опоре}
    \label{fig:p36}
\end{wrapfigure}



% \begin{figure}[H]
%     \centering

%     \begin{tikzpicture}[font=\small,            line cap=round,                             >={Stealth[inset=0pt,round,length=3mm, angle'=20]},vec/.style={>={Stealth[inset=0pt,round,length=3.5mm, angle'=20]}}]


% \filldraw[lightgray,semitransparent] (5,0) rectangle (8,-0.3);
% \draw (5,0) -- (8,0);

% \shade[ball color=lightgray,yshift=0] (6.5,.3) circle (.3);

% %vec
% \node at (6.5,-.2) [circle,fill=NavyBlue,inner sep=0pt,minimum size=1mm,label=below:] {};
% \draw[->,NavyBlue,thick,vec] (6.5,-.2) --++ (0,-1) node[black,right] {$\vec P$};

% \node at (6.5,.2) [circle,fill=red,inner sep=0pt,minimum size=1mm,label=below:] {};
% \draw[->,red,thick,vec] (6.5,.2) --++ (0,1) node[black,right] {$\vec N$};


%     \end{tikzpicture}%
% \caption{Шар на опоре}
% \label{fig:p36}
% \end{figure}

Пусть теперь массивный шар покоится на поверхности планеты~--- то есть, как говорят, на опоре (рис.\,\ref{fig:p36}).

\textbf{Вес} ($\vec P$ [Н])~--- это сила, действующая на
опору или подвес со стороны тела (синий вектор на рис.\,\ref{fig:p36}). \emph{Вес приложен к~опоре (подвесу)}, а не к телу.

\textbf{Сила реакции} ($\vec N$ [Н])~--- это сила, приложенная к телу со стороны опоры или подвеса (красный вектор на рис.\,\ref{fig:p36}).

В рассматриваемой паре тел (шар и опора) силы~$P$ и~$N$ связаны третьим законом Ньютона:
\begin{equation}
    P=N.
\end{equation}

}

Следует отметить, что сила реакции и вес служат проявлением \emph{электромагнитного взаимодействия} тел.






 
\end{document}